% Study guide for learning target 1B. Build: python3 units/ecology/study-guides/build.py 1B
% Macros: latex/preamble.tex (\drawfig, \bookcrop) and latex/guide.sty (\guidetitle, tbpage, \tbnote).
\documentclass[11pt,letterpaper]{article}
% Shared preamble for the study guides and mock exams (built by ../build.py).
% \sgroot is the path from the document's folder to study-guides/ (empty for guides, ../ for exams).
\providecommand{\sgroot}{}
\usepackage[margin=0.8in,headheight=15pt]{geometry}
\usepackage{fontspec}
\setmainfont{texgyrepagella}[Extension=.otf,UprightFont=*-regular,BoldFont=*-bold,ItalicFont=*-italic,BoldItalicFont=*-bolditalic]
\setsansfont{texgyreheros}[Extension=.otf,UprightFont=*-regular,BoldFont=*-bold,ItalicFont=*-italic,BoldItalicFont=*-bolditalic]
\usepackage{amsmath,amssymb,booktabs,longtable,array,calc,enumitem,xcolor,graphicx}
\usepackage{tikz}
\usepackage{pgfplots}\pgfplotsset{compat=1.18}
\usetikzlibrary{calc,arrows.meta,positioning}
\usepackage[most]{tcolorbox}
\usepackage{fancyhdr,titlesec,needspace,etoolbox,environ}
\usepackage{hyperref}

\definecolor{forest}{HTML}{24594B}
\definecolor{pale}{HTML}{EFF5F1}
\definecolor{keyc}{HTML}{C98A0B}   % key idea: amber
\definecolor{vocabc}{HTML}{2F6FA8} % vocabulary: blue
\definecolor{tipc}{HTML}{3F8A3A}   % make it make sense: green
\definecolor{watchc}{HTML}{B8452F} % watch out: red
\definecolor{linkc}{HTML}{7B4C9A}  % connection: purple
\hypersetup{colorlinks=true,linkcolor=forest,urlcolor=vocabc}

\pagestyle{fancy}\fancyhf{}
\renewcommand{\headrulewidth}{0pt}
\fancyfoot[C]{\small\sffamily\thepage}
\setlength{\parindent}{0pt}\setlength{\parskip}{5pt}
\setlength{\emergencystretch}{2em}
\setlist[itemize]{leftmargin=1.3em,itemsep=2pt,topsep=2pt}
\setlist[enumerate]{leftmargin=1.6em,itemsep=2pt,topsep=2pt}
\renewcommand{\arraystretch}{1.2}
\providecommand{\tightlist}{\setlength{\itemsep}{0pt}\setlength{\parskip}{0pt}}
\providecommand{\pandocbounded}[1]{#1}
\def\LTcaptype{none}

\titleformat{\section}{\Large\sffamily\bfseries\color{forest}}{}{0pt}{}[{\color{forest}\titlerule[0.8pt]}]
\titleformat{\subsection}{\large\sffamily\bfseries\color{forest}}{}{0pt}{}
\titleformat{\subsubsection}{\normalsize\sffamily\bfseries}{}{0pt}{}
\titlespacing*{\section}{0pt}{14pt}{8pt}
\titlespacing*{\subsection}{0pt}{12pt}{4pt}
\setcounter{secnumdepth}{0}

% Block quotes (textbook excerpts) as shaded boxes.
\renewenvironment{quote}{\begin{tcolorbox}[breakable,enhanced,colback=pale,colframe=forest,boxrule=0pt,leftrule=2.5pt,arc=0pt,
  left=8pt,right=8pt,top=4pt,bottom=4pt,before skip=4pt,after skip=6pt]\small}{\end{tcolorbox}}

% ------------------------------------------------------------------ annotated textbook pages
% Landscape letter page: the scanned page on the left (full height), notes on the right.
\newlength{\portraitW}\setlength{\portraitW}{8.5in}
\newlength{\portraitH}\setlength{\portraitH}{11in}
\def\pgmargin{21.6}  % pt, top/left margin around the scan
\def\pgheight{568.8} % pt, printed height of the scan (8.5in - 2 margins)
\newcommand{\setpagesize}[2]{%
  \global\pdfpagewidth=#1\global\pdfpageheight=#2%
  \global\paperwidth=#1\global\paperheight=#2}

\newcommand{\circlednum}[2]{\tikz[baseline=(c.base)]\node[circle,fill=#1,text=white,font=\scriptsize\sffamily\bfseries,inner sep=1.2pt,minimum size=11pt](c){#2};}

% \textbookpage{image}{page width pt}{page height pt}{lesson}{book page}{intro}{highlights}{notes}
% Notes (\addnote) are typeset into boxes first, then laid out: each sits level with its highlight
% when there is room, is pushed down to avoid the note above, and is pulled back up from the page bottom.
\newcount\nnotes
\newbox\hdrbox
\newdimen\prevbottom \newdimen\limitdim \newdimen\tmpdim
\def\notegap{7pt}
\newcommand{\addnote}[5]{% color, number, label, anchor (scan pt, <0 = none), text
  \global\advance\nnotes by 1
  \expandafter\ifx\csname nb@\the\nnotes\endcsname\relax\expandafter\newbox\csname nb@\the\nnotes\endcsname\fi
  \global\expandafter\setbox\csname nb@\the\nnotes\endcsname=\hbox{\tikz\node[notebox=#1]{\notehead{#1}{#2}{#3}#5};}%
  \expandafter\xdef\csname na@\the\nnotes\endcsname{#4}%
  \expandafter\xdef\csname nc@\the\nnotes\endcsname{#1}}
\newcommand{\nb}[1]{\csname nb@#1\endcsname}
\newcommand{\layoutnotes}{%
  \global\prevbottom=\dimexpr\pgmargin pt+\ht\hdrbox+\dp\hdrbox\relax
  \ifnum\nnotes>0
    \foreach \k in {1,...,\the\nnotes}{%
      \edef\a{\csname na@\k\endcsname}%
      \tmpdim=\dimexpr\prevbottom+\notegap\relax
      \ifdim \a pt<0pt\else
        \pgfmathsetlength{\dimen0}{\pgmargin+\a*\s-7}%
        \ifdim\dimen0>\tmpdim \tmpdim=\dimen0 \fi
      \fi
      \expandafter\xdef\csname ny@\k\endcsname{\the\tmpdim}%
      \global\prevbottom=\dimexpr\tmpdim+\ht\nb{\k}+\dp\nb{\k}\relax}%
    \global\limitdim=\dimexpr612pt-32pt\relax
    \foreach \k in {\the\nnotes,...,1}{%
      \tmpdim=\csname ny@\k\endcsname
      \ifdim\dimexpr\tmpdim+\ht\nb{\k}+\dp\nb{\k}\relax>\limitdim
        \tmpdim=\dimexpr\limitdim-\ht\nb{\k}-\dp\nb{\k}\relax
        \expandafter\xdef\csname ny@\k\endcsname{\the\tmpdim}%
      \fi
      \global\limitdim=\dimexpr\tmpdim-\notegap\relax}%
    \ifdim\limitdim<\dimexpr\pgmargin pt+\ht\hdrbox+\dp\hdrbox-\notegap\relax
      \typeout{WARNING: notes overflow on textbook page \currentbookpage}\fi
  \fi}
\newcommand{\drawnotes}{%
  \ifnum\nnotes>0
    \foreach \k in {1,...,\the\nnotes}{%
      \edef\a{\csname na@\k\endcsname}\edef\c{\csname nc@\k\endcsname}\edef\y{\csname ny@\k\endcsname}%
      \node[anchor=north west,inner sep=0pt] at ($(O)+(\notex pt,-\y)$) {\usebox{\nb{\k}}};
      \ifdim \a pt<0pt\else
        \draw[\c,line width=0.6pt] ($(scan.north east)+(0,{-\a*\s pt})$) -- ++(6pt,0) -- ($(O)+(\notex pt-1pt,-\y-7pt)$);
      \fi}%
  \fi}
\newcommand{\textbookpage}[8]{%
  \clearpage\setpagesize{11in}{8.5in}\thispagestyle{empty}%
  \def\currentbookpage{#5}%
  \pgfmathsetmacro{\s}{\pgheight/#3}%
  \pgfmathsetmacro{\pgright}{\pgmargin+#2*\s}%
  \pgfmathsetmacro{\notex}{\pgright+20}%
  \pgfmathsetmacro{\notew}{792-\notex-\pgmargin}%
  \pgfmathsetmacro{\notetw}{\notew-10}%
  \setbox\hdrbox=\vbox{\hsize=\notew pt\parskip=0pt
    {\raggedright\sffamily\bfseries\color{forest}#4\par}\vspace{1pt}
    {\sffamily\footnotesize\color{black!55}Miller \& Levine Biology, textbook page #5\par}\vspace{-1pt}
    {\color{forest}\rule{\notew pt}{0.6pt}}\par\vspace{2pt}
    \ifx\relax#6\relax\else{\small\raggedright #6\par}\fi}%
  \global\nnotes=0 #8\layoutnotes
  \null
  \begin{tikzpicture}[remember picture,overlay]
    \coordinate (O) at (current page.north west);
    \node[anchor=north west,inner sep=0pt,draw=black!25] (scan) at ($(O)+(\pgmargin pt,-\pgmargin pt)$)
      {\includegraphics[height=\pgheight pt]{#1}};
    \node[anchor=north west,inner sep=0pt] at ($(O)+(\notex pt,-\pgmargin pt)$) {\usebox{\hdrbox}};
    \node[anchor=south east,font=\scriptsize\sffamily,text=black!50,inner sep=0pt] at ($(O)+(792pt-\pgmargin pt,-603pt)$)
      {Part 1 · the textbook, annotated \quad·\quad Miller \& Levine Biology (2019), p.\,#5 \quad·\quad page \thepage};
    #7
    \drawnotes
  \end{tikzpicture}%
  \clearpage\setpagesize{\portraitW}{\portraitH}}

% Highlight one text line on the scan: \hl{color}{x0}{y0}{x1}{y1} in scan points (y down).
\newcommand{\hl}[5]{%
  \fill[#1,opacity=0.22] ($(scan.north west)+({(#2-1.5)*\s pt},{-(#3-1)*\s pt})$) rectangle
                         ($(scan.north west)+({(#4+1.5)*\s pt},{-(#5+1)*\s pt})$);}
% Numbered badge just left of a highlight.
\newcommand{\badge}[4]{%
  \node[circle,fill=#1,text=white,font=\scriptsize\sffamily\bfseries,inner sep=1pt,minimum size=10pt]
    at ($(scan.north west)+({#3*\s pt-7pt},{-#4*\s pt})$) {#2};}
\newcommand{\notehead}[3]{{\sffamily\bfseries\footnotesize\circlednum{#1}{#2}\ \color{#1}\MakeUppercase{#3}}\\[1pt]}
\tikzset{notebox/.style={text width=\notetw pt,inner sep=4pt,fill=#1!7,draw=#1!55,line width=0.4pt,
  font=\small,align=left,rounded corners=1.5pt}}

% ------------------------------------------------------------------ figures in the guide text
% Caption used under both book figures and new diagrams.
\newcommand{\figcaption}[2]{\par\smallskip{\small\sffamily\raggedright #1\par}%
  \ifx\relax#2\relax\else{\scriptsize\sffamily\color{black!55}Source: #2\par}\fi}
% \bookfig{png}{width fraction}{caption}{source}: a figure cropped from one of the books.
\newcommand{\bookfig}[4]{\par\Needspace{8\baselineskip}\medskip
  \begin{minipage}{\linewidth}\centering
    \fcolorbox{black!15}{white}{\includegraphics[width=#2\linewidth]{#1}}
    \figcaption{#3}{#4}
  \end{minipage}\par\medskip}
% Several book figures side by side: \bookfigrow{\bookfigcol{w}{png}{height}{caption}{source}\hfill ...}
\newcommand{\bookfigcol}[5]{\begin{minipage}[t]{#1\linewidth}\centering
    \fcolorbox{black!15}{white}{\includegraphics[height=#3,width=0.95\linewidth,keepaspectratio]{#2}}
    \figcaption{#4}{#5}\end{minipage}}
\newcommand{\bookfigrow}[1]{\par\Needspace{10\baselineskip}\medskip\noindent#1\par\medskip}
% New diagrams drawn for the guide: \begin{guidefig}...\end{guidefig}{caption}
\newenvironment{guidefig}[1]{\par\Needspace{10\baselineskip}\medskip\begin{minipage}{\linewidth}\centering\def\guidefigcap{#1}}
  {\figcaption{\guidefigcap}{}\end{minipage}\par\medskip}
\tikzset{
  gbox/.style={draw=#1,fill=#1!8,rounded corners=3pt,align=center,font=\small\sffamily,inner sep=5pt},
  garr/.style={-{Stealth[length=2.4mm]},line width=1pt,draw=#1},
  glab/.style={font=\scriptsize\sffamily,fill=white,inner sep=1.5pt,align=center},
}

% ------------------------------------------------------------------ figures shared by guides and exams
% \drawfig{name}: a TikZ/pgfplots figure in figures/name.tex
\newcommand{\drawfig}[1]{\input{\sgroot figures/#1}}
% \bookcrop[width]{id}{book}{page}{x0 y0 x1 y1}{label}{caption}: a figure cropped from a book page.
% book: C = Campbell Biology, ML = Miller & Levine Biology; page = printed page number;
% crop box in points from the page's top-left corner. build.py makes build/figs/<id>.png from these arguments.
\newcommand{\bookname}[1]{\ifstrequal{#1}{C}{Campbell Biology}{Miller \& Levine Biology}}
\newcommand{\bookcrop}[7][0.7]{\bookfig{\sgroot build/figs/#2.png}{#1}{#7}{\bookname{#3}, #6, p.\,#4}}

\usepackage{latex/guide}
\hypersetup{pdftitle={1B Study Guide},pdfauthor={Prepared for Shawn}}
\InputIfFileExists{build/anchors/1B.tex}{}{}

\begin{document}
\guidetitle{1B}{1B Study Guide: What Is Life? Metabolism, Homeostasis, and Feedback}
  {(Textbook 1.3) Scientific investigation includes a variety of methods to develop explanations that are based on evidence and revised through argumentation.}

\section{How to use this guide}

You don't need the textbook to use this guide. Every textbook page that matters for 1B is copied into \textbf{Part 1}, with the important sentences highlighted. Each highlight has a numbered badge that matches a note in the column beside the page. The notes explain hard words, answer the textbook's questions, and connect the page to what your teacher asks about in class.

Part 1 has three groups of pages:

\begin{itemize}
  \item \textbf{Lesson 1.3} (pp. 22–25 and the review on p. 29): what makes something alive, metabolism, homeostasis. This is the lesson the unit guide lists for 1B.
  \item \textbf{Chapter 27} (pp. 907, 908, 932): how feedback keeps the body stable. The class textbook explains feedback here, much later in the book.
  \item \textbf{Lessons 1.1 and 1.2} (pp. 12 and 21): controlled experiments, models, and arguing from evidence, the ``scientific investigation'' half of the 1B target.
\end{itemize}

After Part 1 come the other sources (Campbell covers positive feedback, which the class textbook doesn't), a summary, the vocabulary explained in detail, and full answers to the teacher's study questions. Cover the answers in Part 5 and try each question out loud before you check.

Notes are color-coded:
\notelegend

\section{What this guide covers}

\textbf{Where to read}

\begin{longtable}{@{}>{\raggedright\arraybackslash}p{0.23\linewidth}>{\raggedright\arraybackslash}p{0.28\linewidth}>{\raggedright\arraybackslash}p{0.07\linewidth}>{\raggedright\arraybackslash}p{0.28\linewidth}@{}}
\toprule
Source & Section & Pages & What it covers \\
\midrule
\endhead
Class textbook: \emph{Miller \& Levine Biology} (2019) & \textbf{Lesson 1.3 Patterns of Life} (the lesson the unit guide lists) & 22–25 & characteristics of living things, metabolism, homeostasis, viruses, stability and change, mechanism, system models \\
Class textbook & Lesson 1.1 What Is Science? & 12 & controlled experiments, control groups \\
Class textbook & Lesson 1.2 Science in Context & 21 & models, argument from evidence \\
Class textbook & Lesson 27.1 Organization of the Human Body & 907–908 & homeostasis, negative feedback (feedback inhibition), thermostat, body temperature \\
Class textbook & Lesson 27.3 Human Systems II & 932 & negative feedback in water balance and blood glucose, feedback loop \\
Reference: \emph{Campbell Biology} & Ch. 1, Concept 1.1 & 3, 9 & properties of life, negative and positive feedback \\
Reference: \emph{Campbell Biology} & Ch. 8, Concept 8.1 & 146 & metabolism, catabolic and anabolic pathways \\
Reference: \emph{Campbell Biology} & Ch. 40, Concept 40.2 & 893–894 & homeostasis, set point, dynamic equilibrium, positive feedback in childbirth \\
Teacher materials & What Is Life? video lecture and skeleton notes & — & characteristics of living things, viruses \\
Teacher materials & Feedback Mechanisms and Homeostasis slides; Feedback practice & slides 2–10 & feedback loops, body temperature, childbirth, climate feedbacks \\
\bottomrule
\end{longtable}

Lesson 1.3 introduces homeostasis in one paragraph but doesn't explain feedback. Feedback loops are covered much later in the class textbook (Chapter 27, the human body), and \textbf{positive feedback isn't covered in Miller \& Levine at all}. Those parts use Chapter 27, Campbell, and the teacher's slides.

\clearpage

% ================================================================ Part 1: the textbook, annotated
% Each tbpage is one landscape page: the textbook page on the left, the notes on the right.

\begin{tbpage}{22}{Lesson 1.3 Patterns of Life}
  {The lesson the unit guide lists for 1B. This page answers the first study question: what are the criteria for something to be alive?}
  \tbnote{vocab}{Lesson vocabulary}{sexual reproduction}{evolve}
    {For 1B, the key words are \textbf{metabolism}, \textbf{stimulus}, and \textbf{homeostasis}. Each one is explained on p. 24.}
  \tbnote{key}{}{Biology is the study of life. But what is life?}{}
    {The big question of this lesson: what makes something \textbf{alive}? It's harder to answer than it sounds.}
  \tbnote{watch}{}{No single characteristic is enough}{describe a living thing.}
    {You can't use \textbf{one} trait to decide if something is alive. Living things show \textbf{all} of the characteristics together.}
  \tbnote{tip}{}{a firefly and a fire both give off light}{}
    {A fire gives off light, moves, ``grows,'' and uses energy, but it isn't alive: it has no cells, no DNA, and no homeostasis. That's why you need the whole list.}
  \tbnote{link}{}{such as viruses}{nonliving things.}
    {\textbf{Viruses} are the classic borderline case. The What Is Life? lecture asks whether they're alive. See p. 24 and the answer to study question 1.}
  \tbnote{key}{The 8 criteria for life}{Living things are made up of basic}{changing over time.}
    {\textbf{1} made of cells, \textbf{2} reproduce, \textbf{3} universal genetic code (DNA), \textbf{4} grow and develop, \textbf{5} obtain and use materials and energy, \textbf{6} respond to the environment, \textbf{7} keep a stable internal environment, \textbf{8} evolve. Memorize this list; it answers study question 1.}
\end{tbpage}

\begin{tbpage}{23}{Lesson 1.3 Patterns of Life}
  {The first four characteristics of life.}
  \tbnote{vocab}{organism · cells (criterion 1)}{Living things, or organisms, are made up of}{complex and highly organized.}
    {An \textbf{organism} is an individual living thing. The \textbf{cell} is the smallest unit that is alive. Unicellular organisms are one cell (bacteria); multicellular ones have many (you).}
  \tbnote{vocab}{reproduction (criterion 2)}{All organisms produce new organisms}{has a single parent.}
    {\textbf{Sexual:} two parents; offspring differ from them. \textbf{Asexual:} one parent; offspring are copies. A sterile mule can't reproduce but is still alive, because the criteria describe living things in general, not every individual.}
  \tbnote{vocab}{genetic code (criterion 3)}{Biologists now know that the directions for inheritance}{every organism on Earth.}
    {Instructions for traits are stored in \textbf{DNA}, and nearly every organism uses the same code. That's why it's called \emph{universal}.}
  \tbnote{vocab}{grow and develop (criterion 4)}{All living things grow during at least part}{different functions.}
    {\textbf{Growth} is getting bigger. \textbf{Development} is changing form: one fertilized egg divides into many cells that \textbf{differentiate} into different types (egg $\rightarrow$ caterpillar $\rightarrow$ butterfly).}
\end{tbpage}

\begin{tbpage}{24}{Lesson 1.3 Patterns of Life}
  {The last four characteristics, and the definitions of metabolism, stimulus, and homeostasis.}
  \tbnote{vocab}{metabolism (criterion 5)}{Organisms also need materials and energy just to stay alive.}{life processes is called}
    {\textbf{Metabolism} is \emph{all} the chemical reactions in an organism, both building molecules up (photosynthesis, protein synthesis) and breaking them down (cellular respiration, digestion). This is study question 2.}
  \tbnote{vocab}{stimulus (criterion 6)}{A stimulus is a signal to which an organism responds.}{from within an organism.}
    {A \textbf{stimulus} is a signal an organism reacts to. \textbf{External:} light, temperature (sunflowers turning toward the sun). \textbf{Internal:} from inside the body (thirst). The plural is \emph{stimuli}.}
  \tbnote{vocab}{homeostasis (criterion 7)}{most organisms must keep internal conditions}{you feel thirsty.}
    {\textbf{Homeostasis} means keeping internal conditions (temperature, water) \emph{fairly constant} even when the outside changes. Thirst is an internal stimulus that tells you to drink and restore water balance. It's a \textbf{feedback loop}, explained on pp. 907–908.}
  \tbnote{watch}{evolution (criterion 8)}{As a group, however}{change over time.}
    {\textbf{Groups} evolve over many generations; \textbf{individuals} do not. Your traits don't evolve during your lifetime.}
  \tbnote{link}{Are viruses alive?}{Viruses are particles made}{and evolve.}
    {Viruses \textbf{reproduce and evolve}, but only inside a host cell. They have \textbf{no cells and no metabolism} of their own. Most biologists call them nonliving, or ``at the border.'' Pick a side and support it with the criteria (CER).}
  \tbnote{tip}{Reading check: answer}{Explain Why is it important}{}
    {Cells only work within a narrow range of conditions. If temperature or water levels go too far out of range, cells can't carry out metabolism and the organism can get sick or die (for example, heatstroke or dehydration).}
\end{tbpage}

\begin{tbpage}{25}{Lesson 1.3 Patterns of Life}
  {Crosscutting concepts: big ideas that appear in every unit. Three of them use 1B vocabulary.}
  \tbnote{key}{}{The study of biology revolves around several crosscutting}{function.}
    {Seven big ideas that show up all year. For 1B, focus on \textbf{cause and effect}, \textbf{systems and system models}, and \textbf{stability and change}.}
  \tbnote{vocab}{mechanism}{Cause and Effect: Mechanism and Explanation}{cause and effect.}
    {A \textbf{mechanism} is \emph{how} something happens, the chain of causes and effects. Sweating is a mechanism for cooling; a feedback loop is a mechanism for homeostasis.}
  \tbnote{vocab}{system · model}{Systems and System Models}{circulatory systems.}
    {A \textbf{system} is a set of parts that interact. Blood pressure is controlled by four body systems working together. A \textbf{model} is a simplified picture of a system, like the feedback-loop diagrams you draw in class.}
  \tbnote{key}{Stability and change}{Living things maintain a relatively stable}{fatal consequences.}
    {Homeostasis keeps the body \textbf{stable} while conditions \textbf{change}. If homeostasis breaks down, the result can be serious or fatal, which is why feedback matters (study question 3).}
\end{tbpage}

\begin{tbpage}{29}{Lesson 1.3 Patterns of Life}
  {The end of Lesson 1.3. The top half is lab safety (not on the study questions). The notes answer the lesson review.}
  \tbnote{tip}{Review 1–3: answers}{How do living things differ from nonliving things?}{}
    {\textbf{1.} Living things have \emph{all} the characteristics of life: cells, reproduction, DNA, growth and development, use of materials and energy, response to stimuli, homeostasis, and evolution. Nonliving things may share a few but never all. \textbf{2.} The crosscutting concepts connect all branches of biology with the same big ideas (cause and effect, systems, stability and change, \dots{}). \textbf{3.} Biology can be studied at every scale, from molecules and cells to organisms, ecosystems, and the whole planet.}
  \tbnote{tip}{Review 4–6: answers}{Why do scientists use a common system}{}
    {\textbf{4.} So scientists everywhere can compare and \textbf{replicate} each other's measurements (the metric system). \textbf{5.} All organisms share the same molecules and cell processes, so patterns at that level explain features of all life. \textbf{6.} 250 g $\times$ 10 = 2,500 g = \textbf{2.5 kg}.}
\end{tbpage}

\begin{tbpage}{907}{Lesson 27.1 Organization of the Human Body}
  {From much later in the textbook: how homeostasis actually works. The key idea is negative feedback.}
  \tbnote{vocab}{homeostasis}{This process is called homeostasis}{matter of life or death.}
    {Homeostasis = ``keeping things the \textbf{same}'': relatively constant internal conditions, even when the environment inside or outside changes.}
  \tbnote{vocab}{equilibrium}{Root Words The root word}{}
    {\emph{Stasis} means ``a state of balance or \textbf{equilibrium}.'' Homeo- (same) + stasis (balance) = staying in balance.}
  \tbnote{tip}{dynamic equilibrium}{watched someone driving}{one way or the other.}
    {The car stays in its lane only because the driver keeps making small corrections. That's \textbf{dynamic equilibrium}: stable overall, but always adjusting. The body keeps conditions within a \textbf{range} the same way.}
  \tbnote{vocab}{model · set point}{When the temperature within the house drops below}{within a narrow range.}
    {The thermostat is a \textbf{model} of homeostasis. The \textbf{set point} is the target temperature, and the \textbf{sensor} detects when the room drifts too far from it.}
  \tbnote{vocab}{negative feedback}{Feedback inhibition, or negative feedback, is the process}{opposes the original stimulus.}
    {\textbf{Negative feedback} (feedback inhibition): the response \textbf{opposes} the change. Too cold $\rightarrow$ heat on $\rightarrow$ warmer $\rightarrow$ heat off. It \emph{negates} the change.}
  \tbnote{key}{}{Systems controlled by feedback inhibition}{}
    {Negative feedback makes a system \textbf{stable}. That's why organisms use it for homeostasis.}
  \tbnote{tip}{Reading check: answer}{Describe another example}{}
    {Examples: a car's cruise control (slows $\rightarrow$ more gas $\rightarrow$ speeds up $\rightarrow$ less gas), an air conditioner, an oven, or a toilet tank that refills until the float shuts the valve.}
\end{tbpage}

\begin{tbpage}{908}{Lesson 27.1 Organization of the Human Body}
  {The body's thermostat: negative feedback keeps your temperature near 37\,°C. This is the teacher's slide 4 example.}
  \tbnote{key}{}{Could biological systems achieve homeostasis}{home heating system.}
    {Your body regulates \textbf{temperature} the same way a thermostat does, through negative feedback. Figure 27-4 shows both loops: too hot and too cold.}
  \tbnote{vocab}{sensor · control center}{A part of the brain called the hypothalamus}{monitor body temperature.}
    {The \textbf{hypothalamus} (part of the brain) is the body's thermostat. Its nerve cells \textbf{sense} body temperature and send signals that start or stop warming and cooling.}
  \tbnote{key}{Too cold loop}{If your body temperature drops well below}{body temperature to rise.}
    {Too cold $\rightarrow$ hypothalamus $\rightarrow$ cells speed up \textbf{metabolism} and muscles \textbf{shiver} $\rightarrow$ heat released $\rightarrow$ temperature rises back to normal.}
  \tbnote{key}{Too hot loop}{If body temperature rises too far above}{by evaporation.}
    {Too hot $\rightarrow$ hypothalamus $\rightarrow$ cells slow down and you \textbf{sweat} $\rightarrow$ evaporation cools the skin $\rightarrow$ temperature falls back to normal. Each response opposes the change, so both loops are \textbf{negative feedback}.}
  \tbnote{link}{}{How Do You Respond}{}
    {Blinking when something flies at your face is a \textbf{response to an external stimulus}, another characteristic of life (p. 24).}
\end{tbpage}

\begin{tbpage}{932}{Lesson 27.3 Human Systems II}
  {Two more negative feedback loops: water balance and blood sugar.}
  \tbnote{key}{}{Like most body systems, the endocrine system}{inhibits the system.}
    {Hormone systems are also controlled by \textbf{negative feedback}: a rise in something ``feeds back'' to shut the system down.}
  \tbnote{key}{Water balance loop}{Water balance in the body is one example}{blood water content.}
    {Sweating $\rightarrow$ water lost $\rightarrow$ blood too concentrated $\rightarrow$ hypothalamus $\rightarrow$ ADH (kidneys \textbf{save} water) and \textbf{thirst} (you drink) $\rightarrow$ water back to normal. Drinking a lot $\rightarrow$ less ADH $\rightarrow$ kidneys remove more water. Upper \textbf{and} lower limits = a \textbf{range}.}
  \tbnote{key}{Blood sugar loop}{Glucose concentration in the bloodstream}{back to normal.}
    {After a meal: glucose \textbf{high} $\rightarrow$ pancreas releases \textbf{insulin} $\rightarrow$ cells store glucose $\rightarrow$ level falls. Between meals: glucose \textbf{low} $\rightarrow$ \textbf{glucagon} $\rightarrow$ liver releases glucose $\rightarrow$ level rises. Two opposing hormones keep it in range.}
  \tbnote{vocab}{range}{Insulin and glucagon are opposing hormones}{}
    {``Within a normal \textbf{range}'': homeostasis doesn't keep glucose at one exact number. It keeps it between limits (about 70–110 mg per 100 mL of blood).}
  \tbnote{tip}{Reading check: answer}{Summarize How does the endocrine}{}
    {When blood glucose \textbf{rises}, the pancreas releases \textbf{insulin}, which makes cells take up and store glucose. When it \textbf{falls}, the pancreas releases \textbf{glucagon}, which makes the liver release glucose. Both bring the level back to normal.}
\end{tbpage}

\begin{tbpage}{12}{Lesson 1.1 What Is Science?}
  {Back to scientific investigation, the first half of the 1B target. This page defines the parts of a controlled experiment.}
  \tbnote{vocab}{hypothesis}{Inference, combined with a creative imagination}{support or reject them.}
    {A \textbf{hypothesis} is a testable explanation. Data from experiments either \textbf{support} or \textbf{reject} it; they never ``prove'' it.}
  \tbnote{vocab}{controlled experiment}{Testing hypotheses often involves designing}{controlled experiment.}
    {A \textbf{controlled experiment} changes \textbf{one} variable and keeps everything else the same, so any difference in results must come from that one variable.}
  \tbnote{vocab}{independent · dependent variable}{important to control variables because}{called the dependent variable.}
    {\textbf{Independent variable:} what \emph{you} change on purpose (amount of nitrogen). \textbf{Dependent variable:} what you measure, which \emph{depends} on it (grass height).}
  \tbnote{vocab}{experimental control}{Typically, experiments are divided into control}{a single pair.}
    {The \textbf{control group} gets the same conditions minus the independent variable. It's the comparison that shows whether the variable had any effect. Using several groups (\textbf{replication}) makes results more reliable.}
  \tbnote{link}{A 1B-style experiment}{}{}
    {Does a sports drink help keep body temperature stable during exercise? \textbf{Independent variable:} drink (sports drink or water). \textbf{Dependent variable:} body temperature after 20 minutes of exercise. \textbf{Control group:} the water group. \textbf{Keep constant:} exercise, room temperature, time, and amount drunk.}
\end{tbpage}

\begin{tbpage}{21}{Lesson 1.2 Science in Context}
  {Models, explanations, and arguments: the methods named in the 1B target (``explanations \dots{} revised through argumentation'').}
  \tbnote{vocab}{model}{Developing and Using Models}{communicate ideas to others.}
    {A \textbf{model} can be a drawing, flowchart, physical object, equation, or computer simulation. Your feedback-loop diagrams and the thermostat are both models of homeostasis.}
  \tbnote{key}{}{An explanation describes how variables relate}{supported by evidence.}
    {An \textbf{explanation} tells \emph{how variables relate} and is backed by evidence. Example: ``sweating lowers body temperature because evaporation removes heat.''}
  \tbnote{key}{Argument from evidence}{In science, an argument is a reason}{responding thoughtfully.}
    {This is what the 1B target means by \textbf{revised through argumentation}: scientists use evidence to persuade, listen to critiques, and improve their explanations. CER (claim, evidence, reasoning) is how you practice it.}
  \tbnote{tip}{Review 1–2: answers}{How are both curiosity and skepticism useful}{}
    {\textbf{1.} Curiosity leads scientists to ask new questions; skepticism makes them demand evidence before accepting an answer. \textbf{2.} Peer reviewers find mistakes and weak spots, so scientists can fix their methods and explanations before (and after) publishing.}
  \tbnote{tip}{Review 3–6: answers}{Apply Concepts An advertisement claims}{}
    {\textbf{3.} Science affects society (health, environment, laws), and society's needs shape which questions science studies. \textbf{4.} Testing ideas is how scientists and engineers find out whether an explanation or design actually works. \textbf{5.} Be skeptical: the claim hasn't been checked by independent experts, so there's no reliable evidence yet that the drink works. \textbf{6.} The researcher works for the company that sells the pesticide, so the study could be \textbf{biased}. Independent scientists should repeat it.}
\end{tbpage}

% ================================================================ Part 1 (continued)
\section{Part 1 (continued). Other sources for this target}

The class textbook pages are reproduced in full above. These excerpts cover the rest: the reference book \emph{Campbell Biology} and anything the class textbook leaves out.

\subsection{Campbell Biology}

\textbf{Properties of life} (Ch. 1, p. 3). Campbell says life ``defies a simple definition. We recognize life by what living things do.'' Its Figure 1.2 lists seven properties: \textbf{order, evolutionary adaptation, regulation, reproduction, energy processing, growth and development, and response to the environment.}

\textbf{Feedback regulation} (Ch. 1, p. 9)

\begin{quote}
\textbf{In feedback regulation, the output or product of a process regulates that very process.} The most common form of regulation in living systems is negative feedback, a loop in which the response acts in opposition to the initial stimulus. \dots{} Though less common than processes regulated by negative feedback, there are also many biological processes regulated by \textbf{positive feedback, in which an end product reinforces or strengthens the stimulus.} The clotting of your blood in response to injury is an example.
\end{quote}

\textbf{Metabolism} (Ch. 8, p. 146)

\begin{quote}
\textbf{The totality of an organism's chemical reactions is called metabolism} (from the Greek \emph{metabole}, change). \dots{} Some metabolic pathways release energy by breaking down complex molecules to simpler compounds. These degradative processes are called catabolic pathways, or breakdown pathways. \dots{} Anabolic pathways, in contrast, consume energy to build complicated molecules from simpler ones; they are sometimes called biosynthetic pathways.
\end{quote}

\textbf{Homeostasis and the normal range} (Ch. 40, p. 893)

\begin{quote}
In achieving homeostasis, animals maintain a ``steady state''—a relatively constant internal environment—even when the external environment changes significantly. \dots{} For example, humans maintain a fairly constant body temperature of between approximately 36° and 38°C (97°–99°F), a blood pH within 0.1 pH unit of 7.4, and a blood glucose concentration that is predominantly in the range of 70–110 mg of glucose per 100 mL of blood.
\end{quote}

\textbf{Dynamic equilibrium and positive feedback} (Ch. 40, p. 894)

\begin{quote}
\textbf{Homeostasis is a dynamic equilibrium, an interplay between external factors that tend to change the internal environment and internal control mechanisms that oppose such changes.} \dots{} Unlike negative feedback, positive feedback is a control mechanism that amplifies the stimulus. In animals, positive-feedback loops do not play a major role in homeostasis, but instead help drive processes to completion. During childbirth, for instance, the pressure of the baby's head against sensors near the opening of the mother's uterus stimulates the uterus to contract.
\end{quote}

% ================================================================ Part 2. Summary
\section{Part 2. Summary}

\textbf{The big idea of 1B.} Scientists answer ``What is life?'' and ``How do living things stay alive?'' with evidence from many methods: observation, controlled experiments, and models. They keep revising their explanations when new evidence or better arguments appear. Two ideas sit at the center:

\textbf{1. What counts as alive.} No single trait defines life. Fire moves and gives off light, but it isn't alive. Instead, biologists use a \emph{set of criteria}. According to Miller \& Levine, living things:
\begin{enumerate}
  \item are made of \textbf{cells}
  \item \textbf{reproduce} (sexually or asexually)
  \item are based on a universal \textbf{genetic code} (DNA)
  \item \textbf{grow and develop}
  \item \textbf{obtain and use materials and energy} (metabolism)
  \item \textbf{respond to their environment} (stimuli)
  \item \textbf{maintain a stable internal environment} (homeostasis)
  \item \textbf{evolve}, changing over time as a group
\end{enumerate}

Something is considered alive when it shows \emph{all} of these, not just one or two. \textbf{Viruses} are the classic borderline case: they reproduce and evolve, but only inside a host cell, and they aren't made of cells. Whether to call them alive is a real scientific argument, which is exactly what the 1B target statement describes.

\textbf{2. How living things stay stable.}
\begin{itemize}
  \item \textbf{Metabolism} is all the chemical reactions in an organism. Some build molecules up (photosynthesis, protein synthesis), and some break them down (cellular respiration, digestion).
  \item \textbf{Homeostasis} is keeping internal conditions (temperature, water, blood sugar, pH) within a safe \textbf{range} even when the outside world changes. It is a \textbf{dynamic equilibrium}: conditions constantly drift and get corrected, rather than staying perfectly fixed.
  \item Organisms maintain homeostasis through \textbf{feedback loops}, where a change in a condition triggers a response that affects that same condition.
  \begin{itemize}
    \item \textbf{Negative feedback} \emph{opposes} the change, which keeps the system stable. Examples: sweating and shivering for body temperature, insulin and glucagon for blood sugar. This is how homeostasis works.
    \item \textbf{Positive feedback} \emph{amplifies} the change, pushing the system further in the same direction until something ends it. Examples: childbirth contractions, blood clotting, the melting of ice speeding up global warming. It does \emph{not} maintain homeostasis, and in nature it's often harmful.
  \end{itemize}
\end{itemize}

\textbf{One-sentence version:} life is recognized by a set of shared characteristics, and one of the most important is homeostasis, which organisms maintain mostly through negative feedback loops that keep conditions within a stable range.

\bookcrop[0.85]{c-properties}{C}{3}{60 222 606 752}{Figure 1.2}
  {Some properties of life: order, evolutionary adaptation, regulation (homeostasis), reproduction, energy processing, growth and development, and response to the environment.}

% ================================================================ Part 3. Biology vocabulary
\section{Part 3. Biology vocabulary}

\subsection{organism}
\begin{itemize}
  \item \textbf{Definition:} an individual living thing. ``Living things, or organisms, are made up of small, self-contained units called cells'' (M\&L p. 23).
  \item \textbf{Unicellular vs. multicellular} (from the What Is Life? notes): a unicellular organism is a single cell (bacteria, amoeba). A multicellular organism is made of many cells (a tree, a human).
  \item \textbf{Examples:} a bacterium, a mushroom, a sunflower, a person. \textbf{Non-examples:} a virus (borderline, usually not counted), a cell taken out of your body (alive, but not an organism by itself), a fire.
  \item \textbf{Link to ecology:} organisms are the smallest level in ecology. Organisms $\rightarrow$ population $\rightarrow$ community $\rightarrow$ ecosystem $\rightarrow$ biosphere.
\end{itemize}

\bookcrop[0.8]{c-levels}{C}{4}{28 262 612 752}{Figure 1.3}
  {Levels of biological organization, zooming in from the biosphere to one organism. An organism is an individual living thing; many organisms of one species form a population.}

\subsection{metabolism}
\begin{itemize}
  \item \textbf{Definition:} ``The combination of chemical reactions [an organism uses] as it carries out its life processes'' (M\&L p. 24). Campbell (p. 146): ``The totality of an organism's chemical reactions.'' The teacher's definition: ``All of the chemical reactions that build up and break down molecules in living things.''
  \item \textbf{Two directions:}
\end{itemize}

  | Building up (anabolic) | Breaking down (catabolic) |   |---|---|   | uses energy | releases energy |   | small molecules $\rightarrow$ big molecules | big molecules $\rightarrow$ small molecules |   | photosynthesis (CO\textsubscript{2} + H\textsubscript{2}O $\rightarrow$ glucose) | cellular respiration (glucose $\rightarrow$ CO\textsubscript{2} + H\textsubscript{2}O + energy) |   | protein synthesis (amino acids $\rightarrow$ protein) | digestion (starch $\rightarrow$ sugars) |

\begin{itemize}
  \item \textbf{Memory hook:} \emph{meta} + \emph{bol} comes from the Greek \emph{metabole}, ``change'' (Campbell p. 146). Metabolism is all the chemical \emph{changes} in your body.
  \item \textbf{Why it matters for 1C:} photosynthesis and cellular respiration are how energy and matter move through ecosystems.
\end{itemize}

\drawfig{1B-metabolism}

\subsection{homeostasis}
\begin{itemize}
  \item \textbf{Definition:} ``the relatively constant internal conditions that organisms maintain despite changes in internal and external environments'' (M\&L p. 907). The teacher's note: ``Organisms use feedback mechanisms to maintain internal conditions within certain limits, even as external conditions change.''
  \item \textbf{Word parts:} \emph{homeo-} = same, \emph{-stasis} = ``a state of balance or equilibrium'' (M\&L p. 907). Together: ``keeping things the same.''
  \item \textbf{What gets kept stable in humans} (Campbell p. 893):
  \begin{itemize}
    \item body temperature: about 36–38\,°C (97–99 °F)
    \item blood pH: within 0.1 of 7.4
    \item blood glucose: mostly 70–110 mg per 100 mL
    \item water balance
  \end{itemize}
  \item \textbf{Key point:} homeostasis doesn't mean \emph{no} change. The value moves up and down within a range, which is why it's called a \textbf{dynamic equilibrium} (Campbell p. 894; teacher's slide 2).
  \item \textbf{Why it matters:} ``any breakdown of homeostasis may have serious or even fatal consequences'' (M\&L p. 25). Heatstroke, hypothermia, dehydration, and uncontrolled diabetes are all failures of homeostasis.
\end{itemize}

\drawfig{1B-dynamic-eq}

\subsection{stable / stability}
\begin{itemize}
  \item \textbf{Definition:} \emph{stable} means staying about the same and resisting change. \emph{Stability} is the quality of being stable.
  \item \textbf{In biology:} a stable internal environment is one of the characteristics of life. ``Systems controlled by feedback inhibition are generally very stable'' (M\&L p. 907).
  \item \textbf{Stable $\neq$ unchanging:} a stable body temperature still rises a little during exercise and falls a little at night. It just returns toward the normal value.
  \item \textbf{Crosscutting concept:} ``Stability and Change'' is one of the big themes of biology (M\&L p. 25). Later in this unit you'll see it again: stable ecosystems (2A) and stable populations (2C).
\end{itemize}

\subsection{feedback loop}
\begin{itemize}
  \item \textbf{Definition:} ``any mechanism in which a condition triggers some action that causes a change in that same condition'' (teacher's slide 2). Campbell (p. 9): ``the output or product of a process regulates that very process.''
  \item \textbf{Why it's called a loop:} the result circles back to affect the starting point.
\end{itemize}

  \textbf{condition changes $\rightarrow$ sensor detects it $\rightarrow$ control center $\rightarrow$ response $\rightarrow$ the condition changes back} (and the loop starts again)

\begin{itemize}
  \item \textbf{Parts} (Campbell pp. 893–894; M\&L p. 907):
  \begin{itemize}
    \item \textbf{Stimulus / condition:} the thing being regulated, such as body temperature
    \item \textbf{Sensor:} detects the change, such as nerve cells in the hypothalamus or a thermostat
    \item \textbf{Set point:} the target value, such as 37\,°C
    \item \textbf{Effector / response:} what acts to change the condition, such as sweat glands, muscles (shivering), or a furnace
  \end{itemize}
  \item \textbf{Two types:} negative and positive (below).
\end{itemize}

\bookcrop[0.48]{c-40-8}{C}{894}{25 48 306 592}{Figure 40.8}
  {(a) A thermostat keeps room temperature near its set point. (b) Every negative feedback loop has the same parts: a sensor, a control center, and an effector whose response pushes the variable back toward the set point.}

\subsection{negative feedback}
\begin{itemize}
  \item \textbf{Definition:} ``the process in which a stimulus produces a response that opposes the original stimulus'' (M\&L p. 907). Also called \textbf{feedback inhibition}.
  \item \textbf{What it does:} it \emph{reverses} a change and brings the condition back toward the set point. It keeps things \textbf{stable}, so it's the main way organisms maintain homeostasis. It is ``the most common form of regulation in living systems'' (Campbell p. 9).
  \item \textbf{Examples:}
  \begin{itemize}
    \item \textbf{Body temperature} (teacher's slide 4; M\&L p. 908): too hot $\rightarrow$ sweat $\rightarrow$ cools down. Too cold $\rightarrow$ shiver $\rightarrow$ warms up.
    \item \textbf{Blood glucose} (M\&L p. 932; Campbell p. 9): high after a meal $\rightarrow$ pancreas releases insulin $\rightarrow$ cells take up and store glucose $\rightarrow$ level falls. Low between meals $\rightarrow$ glucagon $\rightarrow$ liver releases glucose $\rightarrow$ level rises.
    \item \textbf{Water balance} (M\&L pp. 24, 932): water is low $\rightarrow$ you feel thirsty and the kidneys save water (ADH) $\rightarrow$ water returns to normal.
    \item \textbf{Thermostat} (M\&L p. 907): too cold $\rightarrow$ furnace on $\rightarrow$ warms $\rightarrow$ furnace off.
    \item \textbf{Climate} (teacher's slide 9): Earth warms $\rightarrow$ more evaporation $\rightarrow$ more clouds $\rightarrow$ more sunlight reflected $\rightarrow$ cooling.
  \end{itemize}
  \item \textbf{Memory hook:} negative = \emph{negates} (cancels) the change.
\end{itemize}

\bookcrop[0.52]{c-glucose-fb}{C}{10}{25 48 306 302}{Figure 1.10}
  {Negative feedback for blood sugar: high glucose $\rightarrow$ the pancreas releases insulin $\rightarrow$ cells take up glucose $\rightarrow$ glucose falls $\rightarrow$ insulin release slows.}

\subsection{positive feedback}
\begin{itemize}
  \item \textbf{Definition:} ``an end product reinforces or strengthens the stimulus'' (Campbell p. 9). ``A control mechanism that amplifies the stimulus'' (Campbell p. 894).
  \item \textbf{What it does:} it pushes the condition \emph{further in the same direction}, so the change keeps growing until something outside the loop stops it. It does \textbf{not} maintain homeostasis. In animals it helps ``drive processes to completion'' (Campbell p. 894). In natural systems it is often harmful (teacher's slide 3).
  \item \textbf{Examples:}
  \begin{itemize}
    \item \textbf{Childbirth} (teacher's slide 5; Campbell p. 894): baby's head presses on the cervix $\rightarrow$ nerve signals to the brain $\rightarrow$ pituitary releases oxytocin $\rightarrow$ stronger contractions $\rightarrow$ more pressure $\rightarrow$ even more oxytocin \dots{} until the baby is born.
    \item \textbf{Blood clotting} (Campbell p. 9): platelets gather at a wound and release chemicals that attract more platelets, until the clot seals the wound.
    \item \textbf{Ice–albedo} (teacher's slide 7): warming melts ice $\rightarrow$ darker surface absorbs more energy $\rightarrow$ more warming $\rightarrow$ more melting.
    \item \textbf{Permafrost} (slide 8) and \textbf{soil decomposition} (slide 10): warming releases more greenhouse gases $\rightarrow$ more warming.
  \end{itemize}
  \item \textbf{Memory hook:} positive = \emph{adds to} the change. ``Positive'' doesn't mean ``good.'' It means ``in the same direction.''
\end{itemize}

\bookcrop[0.55]{c-40-9}{C}{894}{310 48 588 252}{Figure 40.9}
  {(a) Negative feedback: the response reduces the stimulus. (b) Positive feedback: the response increases the stimulus.}

\drawfig{1B-childbirth}

\subsection{experimental control}
\begin{itemize}
  \item \textbf{Definition:} the part of an experiment that gives you something to compare against. It usually means the \textbf{control group}, which ``is exposed to the same conditions as the experimental group except for changes in the independent variable'' (M\&L p. 12). It also refers to keeping all other variables \textbf{controlled} (held constant).
  \item \textbf{Why it matters:} without a control you can't tell whether your independent variable caused the result or something else did. In the salt marsh experiment, the no-nitrogen plots were the control (M\&L p. 11).
  \item \textbf{Example tied to 1B:} to test whether a sports drink helps keep body temperature stable during exercise, one group drinks the sports drink (experimental). The control group drinks plain water, with the same exercise, room temperature, and time. Only the drink differs.
  \item \textbf{Don't confuse:} a \emph{control group} (the comparison group) with \emph{controlled variables} (factors kept the same in both groups). Both are part of experimental control.
\end{itemize}

\drawfig{1B-control-exp}

% ================================================================ Part 4. Academic vocabulary
\section{Part 4. Academic vocabulary}

\subsection{condition}
\begin{itemize}
  \item \textbf{Meaning:} the state something is in, or a factor in the surroundings. \emph{Internal conditions} are inside the body (temperature, water content). \emph{External conditions} are outside (air temperature, sunlight).
  \item \textbf{In 1B:} ``most organisms must keep internal conditions, such as temperature and water content, fairly constant to survive'' (M\&L p. 24). In a feedback diagram, the \textbf{condition} is the box in the middle that's being regulated (teacher's slide 4).
  \item \textbf{In experiments:} ``a control group is exposed to the same \emph{conditions} as the experimental group'' (M\&L p. 12).
\end{itemize}

\subsection{constant}
\begin{itemize}
  \item \textbf{Meaning (adjective):} staying the same, not changing. \textbf{(noun):} something that stays the same, such as ``the constants in an experiment.''
  \item \textbf{In 1B:} homeostasis keeps conditions ``fairly constant'' (M\&L p. 24) or ``relatively constant'' (M\&L p. 907). The words \emph{fairly} and \emph{relatively} matter, because conditions move within a range rather than being perfectly fixed.
  \item \textbf{In experiments:} controlled variables are held \emph{constant}.
\end{itemize}

\subsection{criteria}
\begin{itemize}
  \item \textbf{Meaning:} standards used to judge or decide something. \emph{Criteria} is plural; the singular is \emph{criterion}.
  \item \textbf{In 1B:} ``What are the \textbf{criteria} to determine if something is alive?'' The criteria are the eight characteristics of life. A thing must meet \emph{all} of them to be classified as living.
  \item \textbf{Example sentence:} ``Fire meets some criteria for life, such as using energy and growing, but it does not meet the criteria of being made of cells or having a genetic code.''
\end{itemize}

\subsection{dynamic}
\begin{itemize}
  \item \textbf{Meaning:} always changing or active. The opposite of \emph{static}.
  \item \textbf{In 1B:} homeostasis is a \textbf{dynamic equilibrium}, meaning balance kept by constant adjustment (teacher's slide 2; Campbell p. 894). The textbook also calls the process of science ``dynamic'' because its steps loop back on themselves (M\&L p. 15).
  \item \textbf{Picture it:} a driver keeps ``slightly moving the wheel left and right'' to stay in the lane (M\&L p. 907). The car stays centered, but only because of constant small corrections. That's dynamic stability.
\end{itemize}

\subsection{equilibrium}
\begin{itemize}
  \item \textbf{Meaning:} a state of balance, where opposing forces or processes cancel each other out.
  \item \textbf{In 1B:} the root \emph{stasis} in homeostasis means ``a state of balance or equilibrium'' (M\&L p. 907). In \textbf{dynamic equilibrium}, things keep moving (you eat and excrete every day), but the overall level stays balanced (your weight stays within a few pounds; teacher's slide 2).
  \item \textbf{Static vs. dynamic equilibrium:} a book resting on a table is static equilibrium (nothing moving). Your body temperature is dynamic equilibrium (constant heating and cooling that balance out).
\end{itemize}

\drawfig{1B-static-dynamic}

\subsection{mechanism}
\begin{itemize}
  \item \textbf{Meaning:} the process or set of parts by which something happens; \emph{how} it works.
  \item \textbf{In 1B:} ``The way organisms maintain homeostasis is through \textbf{feedback mechanisms}'' (teacher's slide 2). Shivering and sweating are the mechanisms for body temperature. M\&L's crosscutting concept ``Cause and Effect: \textbf{Mechanism} and Explanation'' (p. 25) means that science explains events by finding the mechanism that causes them.
  \item \textbf{Example sentence:} ``Sweating is a cooling mechanism: evaporation of sweat removes heat from the skin.''
\end{itemize}

\subsection{model}
\begin{itemize}
  \item \textbf{Meaning:} a simplified representation of an object, system, or idea. It can be physical, a drawing, a diagram, a flowchart, a math equation, or a computer simulation (M\&L p. 21).
  \item \textbf{In 1B:}
  \begin{itemize}
    \item The \textbf{feedback loop diagram} you draw in the Feedback practice is a model.
    \item The \textbf{thermostat} is used as a model for homeostasis, ``a nonliving system that keeps conditions within a certain range'' (M\&L p. 907).
  \end{itemize}
  \item \textbf{Why scientists use models:} ``Models help visualize and summarize ideas and communicate ideas to others'' (M\&L p. 21). All models leave something out. The thermostat model, for example, has one sensor, but the body uses many.
\end{itemize}

\bookcrop[0.5]{c-models}{C}{36}{30 48 302 502}{Figure 2.10}
  {Four models of the same molecules: electron diagrams, structural formulas, and space-filling models. Each model shows some features and leaves others out. The thermostat in Figure 40.8 is a model too.}

\subsection{range}
\begin{itemize}
  \item \textbf{Meaning:} the span between a lower limit and an upper limit.
  \item \textbf{In 1B:} homeostasis keeps conditions ``within a certain range, never allowing them to go too far one way or the other'' (M\&L p. 907). Normal human body temperature is about 36–38\,°C; normal blood glucose is 70–110 mg/100 mL (Campbell p. 893).
  \item \textbf{Why it matters:} the teacher's slide 2 says the stable state ``is usually not entirely rigid, but can fluctuate within a certain range.'' Going \emph{outside} the range is when problems start, such as fever, hypothermia, or low blood sugar.
\end{itemize}

\subsection{reinforce}
\begin{itemize}
  \item \textbf{Meaning:} to strengthen or add support to something.
  \item \textbf{In 1B:} positive feedback is when ``an end product \textbf{reinforces} or strengthens the stimulus'' (Campbell p. 9). The Feedback practice asks for ``a \textbf{self-reinforcing} loop,'' a loop that keeps strengthening itself.
  \item \textbf{Example sentence:} ``Melting ice exposes dark water, which absorbs more heat and reinforces the warming.''
  \item \textbf{Everyday use:} practice reinforces learning (a 1A idea).
\end{itemize}

\drawfig{1B-ice-albedo}

% ================================================================ Part 5. Study questions with answers
\section{Part 5. Study questions with answers}

Try each one yourself first, then check.

\subsection{Q1. What are the criteria to determine if something is alive?}

\textbf{Short answer:} No single trait is enough. A living thing shows \emph{all} of these characteristics: it is made of cells, reproduces, has a universal genetic code (DNA), grows and develops, obtains and uses materials and energy, responds to its environment, maintains a stable internal environment (homeostasis), and evolves over time as a group (M\&L p. 22).

\textbf{Each criterion explained}

\begin{longtable}{@{}>{\raggedright\arraybackslash}p{0.04\linewidth}>{\raggedright\arraybackslash}p{0.25\linewidth}>{\raggedright\arraybackslash}p{0.35\linewidth}>{\raggedright\arraybackslash}p{0.24\linewidth}@{}}
\toprule
\# & Criterion & What it means & Example \\
\midrule
\endhead
1 & \textbf{Made of cells} & Cells are ``the smallest units of an organism that can be considered alive'' (M\&L p. 23). Organisms are unicellular (one cell) or multicellular (many cells). & Bacteria (1 cell); humans (trillions) \\
2 & \textbf{Reproduce} & Organisms make new organisms. \textbf{Sexual reproduction:} two parents. \textbf{Asexual reproduction:} one parent (M\&L p. 23). & Budding yeast (asexual); dogs (sexual) \\
3 & \textbf{Universal genetic code} & Inherited instructions are carried in DNA, and the same code works in nearly every organism on Earth (M\&L p. 23). & Offspring resemble their parents \\
4 & \textbf{Grow and develop} & Organisms get bigger, and multicellular organisms \emph{develop}: cells divide and differentiate into different types (M\&L p. 23). & Egg $\rightarrow$ caterpillar $\rightarrow$ monarch butterfly \\
5 & \textbf{Obtain and use materials and energy} & They take in matter and energy and use them in chemical reactions: \textbf{metabolism} (M\&L p. 24). & Plants use sunlight; animals eat food \\
6 & \textbf{Respond to the environment} & They detect and react to \textbf{stimuli}, both external (light, temperature) and internal (thirst) (M\&L p. 24). & Sunflowers turning to follow the sun \\
7 & \textbf{Maintain a stable internal environment} & \textbf{Homeostasis}: internal conditions stay within a safe range (M\&L p. 24). & Sweating when hot \\
8 & \textbf{Evolve} & A \emph{group} of organisms changes over many generations. Individuals don't evolve (M\&L p. 24). & Bacteria evolving antibiotic resistance \\
\bottomrule
\end{longtable}

Campbell's version (p. 3) lists seven properties that overlap closely: order (organized structure), reproduction, growth and development, energy processing, response to the environment, regulation (homeostasis), and evolutionary adaptation.

\bookcrop[0.85]{c-metamorphosis}{C}{725}{60 46 616 352}{Figure 33.42}
  {Growth and development: a butterfly's complete metamorphosis, from larva (caterpillar) to pupa to adult. Development means changing form, not just getting bigger.}

\textbf{Why you need all of them (testing nonliving things):}
\begin{itemize}
  \item \textbf{Fire} grows, uses energy (burns fuel), gives off light, and ``moves in its own way'' (M\&L p. 22). But it has no cells, no DNA, and no homeostasis, so it is \textbf{not alive}.
  \item \textbf{A car} moves, uses energy, and responds (to the driver), but has no cells, doesn't reproduce, and doesn't evolve. \textbf{Not alive.}
  \item \textbf{A crystal} grows, but has no cells or metabolism. \textbf{Not alive.}
  \item \textbf{A mule} is alive even though it's sterile (can't reproduce). It meets every other criterion and is made of cells from parents that did reproduce. The criteria describe living things \emph{in general}, not every individual. (The What Is Life? notes ask: ``Are there any living organisms that cannot reproduce?'')
\end{itemize}

\drawfig{1B-life-checklist}

\textbf{Viruses: the borderline case} (M\&L p. 24; What Is Life? notes)

\begin{longtable}{@{}>{\raggedright\arraybackslash}p{0.38\linewidth}>{\raggedright\arraybackslash}p{0.49\linewidth}@{}}
\toprule
Viruses share with living things & Viruses lack \\
\midrule
\endhead
have genetic material (DNA or RNA) & not made of cells, just proteins, nucleic acids, and sometimes lipids \\
can reproduce, but \emph{only} inside a host cell & have no metabolism of their own \\
evolve (e.g., new flu strains each year) & can't grow or maintain homeostasis \\
can regulate gene expression & can't reproduce on their own \\
\bottomrule
\end{longtable}

Most biologists classify viruses as \textbf{nonliving} (or ``at the border''), because they can't do anything outside a host cell. This question has no single agreed answer, and scientists weigh the evidence and argue both sides. That is exactly the kind of ``explanation revised through argumentation'' the 1B target describes. For a test, state your position and support it with specific criteria, CER style.

\subsection{Q2. Describe and give examples of metabolism.}

\textbf{Definition:} Metabolism is \textbf{all of the chemical reactions that build up and break down molecules in living things} (teacher's note). The class textbook calls it ``the combination of chemical reactions [an organism uses] as it carries out its life processes'' (M\&L p. 24). Campbell calls it ``the totality of an organism's chemical reactions'' (p. 146).

\textbf{Why organisms need it.} Just as ``a building grows taller because workers use energy to assemble new materials'' (M\&L p. 24), organisms need \emph{energy} and \emph{materials} to grow, repair themselves, reproduce, and just stay alive. Metabolism is how they get energy out of food and use it to build what they need.

\textbf{Two kinds of metabolic reactions} (Campbell p. 146):
\begin{enumerate}
  \item \textbf{Breaking down (catabolic):} big, complex molecules are broken into smaller ones, which \textbf{releases energy}.
  \item \textbf{Building up (anabolic):} small molecules are assembled into bigger ones, which \textbf{uses energy}.
\end{enumerate}

Energy released by breakdown reactions powers the building reactions.

\textbf{Examples}

\begin{longtable}{@{}>{\raggedright\arraybackslash}p{0.19\linewidth}>{\raggedright\arraybackslash}p{0.15\linewidth}>{\raggedright\arraybackslash}p{0.33\linewidth}>{\raggedright\arraybackslash}p{0.21\linewidth}@{}}
\toprule
Example & Build up or break down? & What happens & Where \\
\midrule
\endhead
\textbf{Photosynthesis} & Build up & Carbon dioxide + water + light energy $\rightarrow$ glucose + oxygen. Light energy is stored as chemical energy in sugar. & Plants, algae (chloroplasts) \\
\textbf{Cellular respiration} & Break down & Glucose + oxygen $\rightarrow$ carbon dioxide + water + \textbf{energy (ATP)}. This is the main way cells release usable energy. & Nearly all organisms (mitochondria) \\
\textbf{Protein synthesis} & Build up & Amino acids are joined into proteins, following instructions in DNA. & All cells (ribosomes) \\
\textbf{Digestion} & Break down & Large food molecules (starch, proteins, fats) $\rightarrow$ small molecules (sugars, amino acids, fatty acids) the body can absorb. & Digestive system \\
\textbf{Storing glucose as glycogen} & Build up & After a meal, insulin tells liver and muscle cells to link glucose into glycogen (M\&L p. 932). & Liver, muscles \\
\textbf{Breaking down glycogen} & Break down & Between meals, glucagon tells the liver to release glucose (M\&L p. 932). & Liver \\
\bottomrule
\end{longtable}

\bookcrop[0.82]{c-energyflow}{C}{9}{60 48 562 354}{Figure 1.9}
  {Metabolism across a whole ecosystem: photosynthesis builds sugar using light energy, organisms break sugar down to do work, energy leaves as heat, and chemicals cycle back.}

\textbf{Connections:}
\begin{itemize}
  \item \textbf{To homeostasis:} metabolism produces heat. When you're cold, the hypothalamus signals cells to ``speed up their activities,'' and the heat from ``cellular respiration'' raises body temperature (M\&L p. 908). When you're hot, the body slows metabolism down.
  \item \textbf{To ecology (1C):} photosynthesis and cellular respiration are opposites, and together they move energy and carbon through ecosystems. Producers build sugar by photosynthesis; consumers break it down by respiration.
\end{itemize}

\subsection{Q3. How do feedback mechanisms help to maintain homeostasis?}

\textbf{Short answer:} Feedback mechanisms constantly detect changes in internal conditions and trigger responses that push the condition back toward normal. Mostly this happens through \textbf{negative feedback}, which opposes a change. Because every drift away from the set point is corrected, conditions like temperature, water, and blood sugar stay within a safe \textbf{range}, which is homeostasis.

\textbf{1. What homeostasis needs.} Homeostasis means keeping internal conditions ``relatively constant \dots{} despite changes in internal and external environments'' (M\&L p. 907). The outside world keeps changing (hot days, meals, exercise), so the body needs a \emph{control system} that keeps checking and correcting. That system is a \textbf{feedback loop}: ``a condition triggers some action that causes a change in that same condition'' (teacher's slide 2).

\textbf{2. How a negative feedback loop works.} Every loop has the same parts (Campbell pp. 893–894):
\begin{itemize}
  \item a \textbf{sensor} detects the condition
  \item a \textbf{control center} compares it to the \textbf{set point}
  \item an \textbf{effector} produces a \textbf{response} that pushes the condition the opposite way
\end{itemize}

Typically there are two loops working together: ``mechanisms that lower a condition when it is too high, and other mechanisms that raise the level when it is too low'' (teacher's slide 2).

\textbf{3. Example: body temperature} (teacher's slide 4; M\&L p. 908)

\drawfig{temp-loop}

In words:
\begin{itemize}
  \item \textbf{Too hot} (exercise, hot day): nerve cells in the hypothalamus detect it $\rightarrow$ sweat glands produce sweat, and cellular activity slows $\rightarrow$ evaporation cools the skin $\rightarrow$ temperature \textbf{falls} back toward 37\,°C.
  \item \textbf{Too cold:} the hypothalamus detects it $\rightarrow$ cells speed up metabolism, and muscles shiver $\rightarrow$ heat is released $\rightarrow$ temperature \textbf{rises} back toward 37\,°C.
  \item Each response \textbf{opposes} the change that triggered it, which is why it's \emph{negative} feedback. When the temperature is back in range, the response switches off, just like a thermostat turning off a furnace (M\&L p. 907).
\end{itemize}

\bookcrop[0.82]{c-thermo}{C}{901}{60 46 616 430}{Figure 40.18}
  {The hypothalamus works like a thermostat. Too warm: sweating and widened skin blood vessels cool the body. Too cold: shivering and narrowed blood vessels warm it.}

\textbf{4. Example: blood glucose} (M\&L p. 932; Campbell p. 9)
\begin{itemize}
  \item \textbf{After a meal}, blood glucose rises $\rightarrow$ the pancreas releases \textbf{insulin} $\rightarrow$ liver and muscle cells take up glucose and store it as glycogen $\rightarrow$ blood glucose \textbf{falls}.
  \item \textbf{Between meals}, blood glucose falls $\rightarrow$ the pancreas releases \textbf{glucagon} $\rightarrow$ the liver breaks down glycogen and releases glucose $\rightarrow$ blood glucose \textbf{rises}.
  \item Result: blood glucose stays around 70–110 mg per 100 mL (Campbell p. 893).
\end{itemize}

\bookcrop[0.8]{c-glucose}{C}{928}{28 400 614 754}{Figure 41.23}
  {Blood sugar control: insulin lowers blood glucose after eating and glucagon raises it between meals, keeping it near 70–110 mg per 100 mL.}

\textbf{5. Example: water balance} (M\&L pp. 24, 932)
\begin{itemize}
  \item \textbf{Sweating} makes you lose water, so dissolved materials in the blood become more concentrated $\rightarrow$ the hypothalamus detects this $\rightarrow$ it triggers ADH (kidneys save water) and \textbf{thirst} (you drink) $\rightarrow$ water level \textbf{rises} back to normal.
  \item \textbf{After drinking} a lot, less ADH is released $\rightarrow$ the kidneys remove more water $\rightarrow$ the level \textbf{falls} back.
\end{itemize}

\textbf{6. Why it's a \emph{dynamic} equilibrium.} Responses aren't instant, so conditions wobble a little above and below the set point before being corrected. Homeostasis ``moderates but doesn't eliminate changes'' (Campbell p. 894). Stability comes from \emph{constant correction}, like a driver steering to stay in the lane (M\&L p. 907).

\textbf{7. What about positive feedback?} Positive feedback \emph{amplifies} a change instead of reversing it, so it does \textbf{not} maintain homeostasis. In the body it's used only for processes that need to \emph{finish}: childbirth contractions and blood clotting (Campbell pp. 9, 894). In natural systems it's often harmful, because it pushes things to extremes (teacher's slide 3). The climate examples from class:

\begin{longtable}{@{}>{\raggedright\arraybackslash}p{0.07\linewidth}>{\raggedright\arraybackslash}p{0.43\linewidth}>{\raggedright\arraybackslash}p{0.12\linewidth}>{\raggedright\arraybackslash}p{0.24\linewidth}@{}}
\toprule
Slide & Loop & Type & Why \\
\midrule
\endhead
6 & More CO\textsubscript{2} $\rightarrow$ warmer $\rightarrow$ more photosynthesis $\rightarrow$ CO\textsubscript{2} removed $\rightarrow$ cooling & \textbf{Negative} & The result opposes the warming. \\
7 & Warmer $\rightarrow$ ice melts $\rightarrow$ lower albedo $\rightarrow$ more heat absorbed $\rightarrow$ warmer & \textbf{Positive} & The result adds to the warming. \\
8 & Warmer $\rightarrow$ permafrost melts $\rightarrow$ methane released $\rightarrow$ warmer & \textbf{Positive} & The result adds to the warming. \\
9 & Warmer $\rightarrow$ more evaporation $\rightarrow$ more clouds $\rightarrow$ more reflection $\rightarrow$ cooler & \textbf{Negative} & The result opposes the warming. \\
10 & Warmer $\rightarrow$ faster decomposition $\rightarrow$ more CO\textsubscript{2} $\rightarrow$ warmer & \textbf{Positive} & The result adds to the warming. \\
\bottomrule
\end{longtable}

\textbf{How to tell them apart:} follow the loop all the way around. If you end up pushing the original condition \emph{back}, it's negative. If you push it \emph{further the same way}, it's positive.

\textbf{8. When feedback fails.} ``Any breakdown of homeostasis may have serious or even fatal consequences'' (M\&L p. 25). If sweating can't keep up on a very hot day, body temperature climbs out of range (heatstroke). If the insulin loop doesn't work, blood glucose stays high (diabetes).

\textbf{Tips for the Feedback practice worksheet}
\begin{itemize}
  \item For \textbf{negative feedback}, put a \emph{measurable} condition in the center box and draw two loops, one for ``too high'' and one for ``too low.'' Everyday examples: water temperature in a shower (too hot $\rightarrow$ turn the cold up; too cold $\rightarrow$ turn the hot up), or hunger level (hungry $\rightarrow$ eat $\rightarrow$ full $\rightarrow$ stop eating).
  \item For \textbf{positive feedback}, make sure every arrow makes the condition \emph{more extreme}. Everyday examples: a microphone near a speaker (sound $\rightarrow$ amplified $\rightarrow$ picked up again $\rightarrow$ louder squeal), or a rumor spreading (more people know $\rightarrow$ more people tell others $\rightarrow$ even more people know).
  \item Don't use technical examples you found online, such as blood glucose. The worksheet asks for your own everyday examples.
\end{itemize}

% ================================================================ Sources
\section{Sources}

\begin{itemize}
  \item Miller, K. R., \& Levine, J. S. \emph{Miller \& Levine Biology}. Pearson, 2019. Lesson 1.1 (p. 12), Lesson 1.2 (pp. 15, 21), Lesson 1.3 (pp. 22–25), Lesson 27.1 (pp. 907–908), Lesson 27.3 (p. 932).
  \item Urry, L. A., et al. \emph{Campbell Biology}. Ch. 1, Concept 1.1 (pp. 3, 9); Ch. 8, Concept 8.1 (p. 146); Ch. 40, Concept 40.2 (pp. 893–894).
  \item Teacher materials: \emph{Ecology Unit Guide (KJ)}, \emph{What Is Life? Skeleton Notes}, \emph{Feedback Mechanisms and Homeostasis} slides, \emph{Feedback Mechanisms and Homeostasis Practice}.
\end{itemize}

\end{document}
